% =====================================================================
% CHƯƠNG 5. KẾT LUẬN
% Không cần thêm gói nào ngoài preamble của Format_KLTN/main.tex.
% =====================================================================
\chapter{Kết luận}
\label{chap:conclusion}

Khóa luận đã nghiên cứu bài toán tự động cập nhật đồ thị tri thức từ bản tin bằng mô hình ngôn ngữ lớn.
Xuất phát điểm là quan sát rằng KG nhanh chóng lỗi thời so với văn bản,
trong khi các hệ xây dựng KG từ văn bản hiện có chủ yếu chỉ thêm tri thức
và hầu như không vô hiệu hóa được tri thức đã hết hiệu lực \cite{zaporojets2025emerge}.
Chương này tổng kết các kết quả đạt được, rút ra ý nghĩa của chúng,
nêu các hạn chế còn tồn tại và đề xuất hướng phát triển tiếp theo.

% ---------------------------------------------------------------------
\section{Kết quả đạt được}
\label{sec:conclusion-results}

Về phương pháp, khóa luận đã thiết kế và cài đặt một pipeline sáu bước cho bài toán cập nhật KG bi-temporal,
vận hành theo nguyên tắc LLM trích xuất, cấu trúc khoanh vùng và bộ phán quyết định.
Pipeline gồm bước rút câu đơn, bước trích thao tác tường minh với đầu ra JSON ép schema,
bước áp thao tác lên KG theo nguyên tắc không xóa mà chỉ vô hiệu hóa có dấu thời gian,
bước khoanh vùng cấu trúc hai vế không dùng luật miền, bộ phán LLM gom lô và bước gộp tên định kỳ.
Toàn bộ hệ thống chạy với Gemma~4 E4B lượng tử hóa 4 bit \cite{gemma4} qua llama.cpp \cite{llamacpp} trên GPU T4.

Về thực nghiệm, khóa luận đã đánh giá module trích triple theo schema,
bước nền của mọi thao tác thêm, trên ba benchmark WebNLG, REBEL và Wiki-NRE \cite{zhang2024edc,cabot2021rebel}.
EDC \cite{zhang2024edc} và ATOM+EDC \cite{lairgi2025atom} được chạy lại với cùng LLM
để tách đóng góp của phương pháp khỏi đóng góp của mô hình,
và mọi hệ được chấm bằng cùng script đánh giá của EDC.
Với một lệnh gọi LLM mỗi câu, module đề xuất đạt F1 theo chế độ Partial 74,4 trên WebNLG, 55,2 trên REBEL và 66,2 trên Wiki-NRE,
cao hơn EDC chạy cùng LLM (70,6, 49,2 và 62,2) lần lượt 3,8, 6,0 và 4,0 điểm,
và cao hơn rõ rệt ATOM+EDC (56,8, 40,3 và 40,7).
Thời gian chạy của module là 20, 21 và 15 phút, so với 50, 61 và 41 phút của EDC.
Các điểm số này cũng tương đương, thậm chí nhỉnh hơn, số công bố của EDC khi dùng Mistral-7B ở bước trích mở (72,8, 51,6 và 64,7),
dù hai kết quả không so sánh trực tiếp được vì cấu hình công bố dùng GPT-3.5 cho các bước còn lại.

Từ các kết quả trên, khóa luận rút ra ba nhận định.
Thứ nhất, với một LLM mở cỡ nhỏ, việc ép quan hệ vào schema ngay lúc sinh bằng giải mã có ràng buộc
cho kết quả tốt hơn quy trình trích mở rồi mới chuẩn hóa về schema như EDC,
có lẽ vì mô hình không phải tự đặt tên quan hệ rồi trông vào một bước sau để sửa lại.
Thứ hai, thêm bước LLM phía trước như cách ATOM tách văn bản thành mệnh đề không làm tăng mà còn làm giảm F1 trên các câu ngắn;
với mô hình nhỏ, mỗi bước LLM thêm vào là một chỗ có thể làm mất thông tin.
Thứ ba, một lệnh gọi mỗi câu giữ chi phí đủ thấp để triển khai trên phần cứng phổ thông.
Các nhận định này ủng hộ lựa chọn thiết kế của pipeline:
giao cho LLM những tác vụ hẹp với đầu ra ép schema,
và chuyển phần lập luận trên cấu trúc sang các bước tất định.
Phương pháp vì vậy có triển vọng ứng dụng để duy trì KG chuyên ngành, KG nội bộ của tổ chức
hoặc bộ nhớ dài hạn của tác tử LLM, nơi dữ liệu cần được xử lý tại chỗ và tri thức cần được cập nhật liên tục.

% ---------------------------------------------------------------------
\section{Hạn chế}
\label{sec:conclusion-limits}

\begin{itemize}
  \item \textbf{Phần lớn pipeline chưa được đánh giá định lượng.}
  Các bước khoanh vùng cấu trúc, bộ phán, rút câu đơn và gộp tên mới được thiết kế và cài đặt,
  chưa được đánh giá định lượng trong khóa luận.
  Vì vậy, khả năng vô hiệu hóa fact và phát hiện thay đổi gián tiếp,
  vốn là điểm khác biệt chính của phương pháp, mới được lập luận và minh họa qua ví dụ.

  \item \textbf{Chỉ dùng một LLM.}
  Mọi thực nghiệm dùng Gemma~4 E4B lượng tử hóa 4 bit;
  chưa rõ ưu thế so với EDC có còn giữ nguyên với các LLM khác, đặc biệt là các mô hình lớn hơn.

  \item \textbf{So sánh với số công bố chỉ mang tính tham khảo.}
  Mẫu 1.000 câu của REBEL và Wiki-NRE được lấy từ bản dữ liệu trong kho mã của EDC,
  không chắc trùng với mẫu mà bài báo EDC dùng;
  cấu hình công bố của EDC cũng chỉ đổi LLM ở bước trích mở.

  \item \textbf{Benchmark chưa có yếu tố thời gian.}
  Ba benchmark gồm các câu độc lập, không có KG trước và không có mốc thời gian,
  nên chưa đo được việc quy đổi thời gian tương đối và việc nhận ra fact cũ đã hết hiệu lực.
\end{itemize}

% ---------------------------------------------------------------------
\section{Hướng phát triển}
\label{sec:conclusion-future}

\begin{itemize}
  \item \textbf{Đánh giá trên benchmark cập nhật KG:}
  đo các bước còn lại trên EMERGE \cite{zaporojets2025emerge} cho thao tác giữ, thêm và vô hiệu hóa,
  và trên NBAtransactions \cite{tang2019gupdate} cho thay đổi gián tiếp, so với GUpdater;
  đồng thời lặp lại thực nghiệm với nhiều LLM khác nhau.

  \item \textbf{Tạo thực thể mới (Mint):}
  bổ sung bước quyết định khi nào một tên trong bản tin là thực thể chưa có trong KG,
  cấp định danh mới và sinh các fact đi kèm.

  \item \textbf{Bộ phán lai:}
  kết hợp LLM với mô hình dự đoán liên kết quy nạp (\textit{inductive}) như AnyBURL \cite{meilicke2019anyburl} hoặc ULTRA \cite{galkin2024ultra};
  ứng viên có độ tin cậy cao được quyết định trực tiếp, chỉ vùng xám mới chuyển cho LLM, nhằm giảm số lệnh gọi.
  Các mô hình nhúng kiểu \textit{transductive} như TransE \cite{bordes2013transe} hay RotatE \cite{sun2019rotate}
  không phù hợp vì chỉ biểu diễn được các thực thể đã thấy khi huấn luyện.

  \item \textbf{Bộ dữ liệu tiếng Việt:}
  xây dựng một bộ bản tin tiếng Việt kèm KG trước và sau sự kiện để đánh giá phương pháp trên tiếng Việt.

  \item \textbf{Hệ thống trình diễn:}
  xây dựng giao diện cho phép xem lịch sử của từng fact trong KG bi-temporal,
  duyệt và chấp nhận các thao tác do hệ thống đề xuất.
\end{itemize}

Tóm lại, khóa luận đã đề xuất một pipeline cập nhật KG có thao tác vô hiệu hóa và lưu được lịch sử,
và đã cho thấy bước nền của pipeline là trích triple theo schema,
khi chạy với một LLM mở cỡ nhỏ, chính xác hơn EDC và ATOM+EDC, đồng thời nhanh hơn EDC khi dùng cùng mô hình.
Đánh giá định lượng các bước còn lại trên benchmark cập nhật KG là việc cần làm tiếp theo
để đưa phương pháp tới gần một hệ thống cập nhật KG tự động có thể triển khai trong thực tế.
